% !TeX program = xelatex
\documentclass[tikz,border=6pt]{standalone}
\usepackage[UTF8]{ctex}
\usetikzlibrary{arrows.meta,positioning,fit,backgrounds,calc}

\definecolor{hitblue}{RGB}{39,82,128}
\definecolor{inputfill}{RGB}{235,242,248}
\definecolor{detectfill}{RGB}{225,239,247}
\definecolor{selectfill}{RGB}{231,244,240}
\definecolor{evalfill}{RGB}{246,240,225}
\definecolor{improvefill}{RGB}{239,234,246}
\definecolor{goldfill}{RGB}{248,238,232}
\definecolor{textgray}{RGB}{45,50,56}

\tikzset{
  stage/.style={
    rectangle,
    rounded corners=2pt,
    draw=hitblue,
    line width=0.8pt,
    minimum width=3.15cm,
    minimum height=2.35cm,
    align=center,
    inner sep=5pt,
    text=textgray,
    font=\small
  },
  stage title/.style={
    font=\bfseries\small,
    text=hitblue
  },
  flow/.style={
    -{Latex[length=2.5mm,width=1.7mm]},
    line width=0.9pt,
    draw=hitblue
  },
  gold/.style={
    rectangle,
    rounded corners=2pt,
    draw=black!55,
    fill=goldfill,
    line width=0.7pt,
    minimum width=3.15cm,
    minimum height=0.95cm,
    align=center,
    font=\small,
    text=textgray
  },
  goldflow/.style={
    -{Latex[length=2.3mm,width=1.6mm]},
    densely dashed,
    line width=0.8pt,
    draw=black!60
  },
  note/.style={
    font=\scriptsize,
    text=black!65,
    align=center
  }
}

\begin{document}
\begin{tikzpicture}[node distance=0.56cm]

\node[stage,fill=inputfill] (input) {
  {\bfseries\color{hitblue}实验输入}\\[3pt]
  PR 标题与正文\\
  base 版本仓库\\[2pt]
  {\scriptsize 统一任务描述}
};

\node[stage,fill=detectfill,right=of input] (detect) {
  {\bfseries\color{hitblue}测试目标检测}\\[3pt]
  OpenSpec / Spec Kit\\
  cc-sdd / IssueExec\\[2pt]
  {\scriptsize 生成原始目标与理由}
};

\node[stage,fill=selectfill,right=of detect] (select) {
  {\bfseries\color{hitblue}文件级目标选择}\\[3pt]
  路径解析与规范化\\
  存在性、测试属性检查\\
  去重形成候选集合
};

\node[stage,fill=evalfill,right=of select] (evaluate) {
  {\bfseries\color{hitblue}效果评价}\\[3pt]
  命中率与召回率\\
  micro / macro 指标\\
  运行状态与案例分析
};

\node[stage,fill=improvefill,right=of evaluate] (improve) {
  {\bfseries\color{hitblue}方法改进}\\[3pt]
  分析误检与漏检\\
  优化需求表示\\
  关联分析与候选选择
};

\draw[flow] (input) -- (detect);
\draw[flow] (detect) -- (select);
\draw[flow] (select) -- (evaluate);
\draw[flow] (evaluate) -- (improve);

\node[gold,below=1.05cm of detect] (head) {
  head 中实际修改的测试文件
};
\node[gold,right=0.56cm of head] (goldset) {
  独立封存的 gold 集合
};

\draw[goldflow] (head) -- (goldset);
\draw[goldflow] (goldset.east) -| node[pos=0.72,below,note] {仅在评价阶段解封} (evaluate.south);

\coordinate (inputdetectmid) at ($(input.east)!0.5!(detect.west)$);
\coordinate (feedbackbottomright) at ($(improve.south)+(0,-2.75)$);
\coordinate (feedbackbottomleft) at (inputdetectmid |- feedbackbottomright);

% 反馈回路：全程使用带箭头的虚线，终点位于第一条主流程蓝色箭头的中点。
\draw[goldflow]
  (improve.south)
  -- (feedbackbottomright)
  -- node[midway,below=2pt,note] {反馈优化}
     (feedbackbottomleft)
  -- (inputdetectmid);

\begin{scope}[on background layer]
  \node[
    draw=black!35,
    densely dashed,
    rounded corners=3pt,
    fit=(head)(goldset),
    inner sep=7pt,
    label={[note]below:参考答案隔离区}
  ] {};
\end{scope}

\node[note,above=0.16cm of detect] {相同模型、预算、权限与运行环境};
\node[note,below=0.16cm of select] {统一输出为仓库相对测试文件路径};

\end{tikzpicture}
\end{document}
